\begin{lemma}{Gleichmäßige Lipschitzstetigkeit}
             {GleichmaessigeLipschitzstetigkeit}
Seien $(M,d)$ ein metrischer Raum und $I$ eine nichtleere Menge.\\
Sei $(f_i)_{i \in I} \in \mcf(M, \R)^I$ mit den Eigenschaften:
\begin{enumerate}[(a)]
\item $\forall\, m \in M: \menge{f_i(m)}{i \in I}$ ist nach oben beschr�nkt,
\item $\exists\, c \in \R_{>0}\,\forall\, i \in I \, \forall\, x, y \in M:
\abs{f_i(x)-f_i(y)}\le c\cdot d(x, y)$.
\end{enumerate}
Dann ist die Abbildung
\[f: M \rightarrow \R, m \mapsto \sup\menge{f_i(m)}{i \in I}\]
wohldefiniert und lipschitzstetig mit Lipschitzkonstante $c$.
\end{lemma}
\begin{proof}
Die Wohldefiniertheit ist offensichtlich, denn wegen $I \neq \emptyset$ ist
für jedes $m \in M$ die Menge $\menge{f_i(m)}{i \in I}$ nichtleer und wegen
$(a)$ nach oben beschränkt, besitzt also ein Supremum (und auch genau eins).\\
Seien nun $x, y \in M$. O. B. d. A.: $f(x) \ge f(y)$.\\
Dann gilt:
\begin{align*}
\abs{f(x)-f(y)} &= f(x) - f(y) = \left(\sup_{i \in I}f_i(x)\right) -
 \left(\sup_{i \in I}f_i(y)\right)\\
 & = \left(\sup_{i \in I}(f_i(x)-f_i(y) + f_i(y))\right) -
 \left(\sup_{i \in I}f_i(y)\right)\\
 &\le \left(\sup_{i \in I}(f_i(x) -f_i(y)) \right)
 + \left(\sup_{i\in I}f_i(y)\right) - \left(\sup_{i\in I}f_i(y)\right)\\
 & = \sup_{i \in I}\underbrace{(f_i(x) -f_i(y))}_{\le c\cdot d(x, y)}
 \le c \cdot d(x, y)\text.
\end{align*}
Also ist $f$ lipschitzstetig mit Lipschitzkonstante $c$.
\end{proof}